\section{Introdução}

Uma suíte de testes pode executar quase todas as linhas de um sistema e, ainda assim, não perceber um comportamento incorreto. A cobertura de código mede o que foi executado, não o que foi verificado. Em APIs \textit{backend}, esse descompasso pesa porque as regras de negócio combinam limites numéricos, estados e interações entre entidades, e testes baseados em exemplos verificam apenas os cenários que o desenvolvedor lembrou de escrever.

Testes Baseados em Propriedades (PBT, do inglês \textit{Property-Based Testing}) atacam esse problema por outro ângulo: o desenvolvedor descreve uma propriedade geral do sistema e a ferramenta gera automaticamente centenas de entradas para tentar refutá-la \citep{claessen2000quickcheck}. Para medir se uma suíte de fato detecta defeitos, o \MT{} introduz pequenas alterações artificiais no código (mutantes) e verifica quantas delas fazem algum teste falhar; essa proporção é o \ms{} \citep{jia2011mutation,papadakis2019mutation}.

A evidência disponível sobre a efetividade do PBT vem quase sempre de um único ecossistema. \citet{ravi2025pbt} mediram por mutação as suítes de 40 projetos Python e mostraram que testes de propriedade encontram muito mais mutantes que testes unitários; \citet{goldstein2024pbt} e \citet{corgozinho2023pbt} estudaram como desenvolvedores escrevem propriedades, também dentro de um ecossistema. Ao mesmo tempo, ferramentas de mutação diferentes produzem conjuntos de mutantes diferentes para o mesmo programa \citep{kintis2018tools,monteiro2022tools}. Não se sabe, portanto, se as mesmas propriedades, aplicadas a implementações equivalentes em ecossistemas distintos e avaliadas pelas ferramentas de mutação de cada um, levam ao mesmo resultado. Essa é a lacuna deste trabalho.

O estudo compara Java/Spring Boot e Node.js/Express, dois dos ecossistemas mais usados para construir APIs REST, por meio de duas implementações equivalentes de uma API de gestão de eventos e reservas de ingressos. As mesmas oito propriedades de negócio foram escritas com jqwik (Java) e fast-check (Node.js) e avaliadas com PIT e StrykerJS, respectivamente.

\textbf{Questão de pesquisa geral:} suítes de testes baseadas em propriedades, derivadas das mesmas regras de negócio, têm a mesma efetividade, medida por \ms{}, quando aplicadas a implementações equivalentes de uma API REST em Java/Spring Boot e em Node.js/Express?

\textbf{Objetivo geral:} comparar, por meio do \ms{}, a efetividade de suítes PBT equivalentes nas duas implementações. Os objetivos específicos, cada um ligado às questões de pesquisa (QP) da Seção~\ref{sec:gqm}, são:

\begin{enumerate}[label=\textbf{OE\arabic*}, leftmargin=1.6cm]
  \item caracterizar as propriedades derivadas das regras de negócio e sua equivalência entre jqwik e fast-check (QP1);
  \item comparar o \ms{} das duas suítes, inclusive sob o mesmo conjunto de categorias de operadores (QP2--QP4);
  \item analisar a relação entre cobertura e \ms{} em cada ecossistema (QP5);
  \item classificar os mutantes sobreviventes segundo a regra de negócio afetada (QP6);
  \item comparar o custo de execução e o esforço de configuração das ferramentas (QP7).
\end{enumerate}

As contribuições são: (i) um experimento comparativo com 30 execuções por ecossistema e teste estatístico; (ii) a caracterização dos mutantes que sobrevivem a propriedades de invariante em cada ecossistema; e (iii) um repositório reprodutível com a especificação, as duas APIs, as suítes, os scripts de coleta e os dados brutos.

O restante do artigo está organizado assim: a Seção~2 apresenta os conceitos; a Seção~3, os trabalhos relacionados; a Seção~4, o método; a Seção~5, os resultados; a Seção~6, as ameaças à validade; e a Seção~7, as conclusões.
